% !TeX program = lualatex
% Compile twice: lualatex open_problems_500.tex
% All references and prose are embedded; no BibTeX database or external inputs.
\documentclass[11pt,a4paper]{article}
\usepackage[margin=24mm,headheight=14pt]{geometry}
\usepackage{fontspec}
\setmainfont{Latin Modern Roman}
\setsansfont{Latin Modern Sans}
\setmonofont{Latin Modern Mono}
\usepackage{amsmath,amssymb,mathtools}
\usepackage{unicode-math}
\setmathfont{Latin Modern Math}
\usepackage{microtype}
\usepackage{enumitem,xurl,needspace}
\usepackage[dvipsnames]{xcolor}
\usepackage{hyperref}
\hypersetup{unicode=true,colorlinks=true,linkcolor=MidnightBlue,urlcolor=MidnightBlue,
 pdftitle={500 Mathematical Problem Statements},pdfauthor={Source-annotated compilation},bookmarksnumbered=true}
\usepackage{bookmark}
\usepackage{tocloft}
\setlength{\cftsecnumwidth}{3em}
\cftsetpnumwidth{2.5em}
\cftsetrmarg{4em}
\usepackage{titlesec}
\titleformat{\section}{\Large\bfseries\raggedright}{\thesection}{0.7em}{}
\usepackage{fancyhdr}
\pagestyle{fancy}\fancyhf{}
\fancyhead[L]{\small\sffamily Mathematical problem statements}
\fancyhead[R]{\small Source edition 22}
\fancyfoot[C]{\thepage}
\setlength{\parindent}{0pt}
\setlength{\parskip}{5pt plus 1pt}
\setlength{\emergencystretch}{3em}
\setcounter{tocdepth}{1}
\setcounter{secnumdepth}{1}
\setlist[itemize]{leftmargin=3em,itemsep=5pt,topsep=4pt,parsep=0pt}
\allowdisplaybreaks
\newcommand{\N}{\mathbb{N}}
\newcommand{\Z}{\mathbb{Z}}
\newcommand{\Q}{\mathbb{Q}}
\newcommand{\R}{\mathbb{R}}
\newcommand{\C}{\mathbb{C}}
\newcommand{\F}{\mathbb{F}}
\newcommand{\A}{\mathbb{A}}
\newcommand{\PP}{\mathbb P}
\newcommand{\E}{\mathbb{E}}
\newcommand{\eps}{\varepsilon}
\newcommand{\dd}{\,\mathrm d}
\newcommand{\sref}[2]{\hyperlink{source:#1:#2}{[S#2]}}
\newcommand{\eref}[2]{\hyperlink{extra:#1:#2}{[E#2]}}
% Turn the next switch off for a shorter reading copy. The quotations preserve
% every exactTarget field, including qualifications omitted from a short summary.
\newif\ifshowcatalogue
\showcataloguetrue
\newcommand{\cataloguescope}[1]{\ifshowcatalogue
 \paragraph{Exact scope in the supplied catalogue.}
 \begin{quote}\small #1\end{quote}\fi}

% Standalone research notebook, original catalogue problem 186.
% Compile twice with: lualatex 186.tex
\numberwithin{equation}{section}
\hypersetup{pdftitle={Research attempt -- problem 186},pdfauthor={Research notebook; original compilation retained}}
\fancyhead[L]{\small\sffamily Research notebook}
\fancyhead[R]{\small\sffamily Problem 186 | 27 September 2026}
\begin{document}
\setcounter{section}{185}
% ================================================================
% problemId: problem.beilinson-soule-vanishing-conjecture-for-rational-motivic-cohomology
% source releaseRank: 186; source releaseStatus: open_with_solved_subcases
% exactTarget SHA-256: c9ab843ce720a16c99765847804cb6cd459960d41972753c203980ac8e77b789
\clearpage
\section{\texorpdfstring{Beilinson-Soule Vanishing Conjecture for Rational Motivic Cohomology}{Beilinson-Soule Vanishing Conjecture for Rational Motivic Cohomology}}\label{186}
\subsection{Definitions and mathematical statement}
For a smooth scheme $X$ over a field, the selected rational motivic groups are
\[
 H^{n,i}(X)=\operatorname{gr}_{\gamma}^{i}K_{2i-n}(X)\otimes\Q.
\]
Here $K_j$ is algebraic $K$-theory and the gamma filtration is the filtration defined by the lambda-ring operations, with its rational graded pieces providing the weight decomposition. The Beilinson--Soul\'e assertion is
$H^{n,i}(X)=0$ whenever $n\le0$ and $(n,i)\ne(0,0)$, with the conventional zero groups outside the available $K$-theoretic weights. The exceptional group $H^{0,0}$ is not required to vanish. This is a rational statement about all smooth schemes over fields, not the stronger claim that the corresponding integral groups have no torsion.
\par\smallskip\textit{Formulation and concept references:} \sref{186}{1}.
\cataloguescope{Prove that for every smooth scheme X over a field k the rational motivic cohomology groups H\textasciicircum{}(n,i)(X), defined as the i-th graded piece for the gamma filtration on K\_(2i-n)(X) tensored with the rationals, vanish whenever n <= 0 and the pair (n,i) is not (0,0).}
\subsection{Short English statement}
Do all nonpositive-degree rational motivic cohomology groups vanish except for the ordinary degree-zero, weight-zero part?
% BEGIN RESEARCH ADDITION ATTEMPT-186
\subsection{Research attempt}

\paragraph{Current frontier.}
The cited manuscript separates strong Beilinson--Soul\'e vanishing (including degree zero at positive weight) from the weak negative-degree version, and relates the strong version to the motivic $t$-structure for Tate motives \sref{186}{1}, \eref{186}{1}, Section~5.2. These are not general vanishing theorems for arbitrary fields. The attempt below uses the established localization, continuity, and transfer formalism, not existence of the conjectural $t$-structure. The foundational weight-one computation and coniveau machinery provide useful exact tests \eref{186}{2}.

\paragraph{Attack route.}
A direct attack on all smooth schemes seems to introduce geometric complexity before the field obstruction is understood. Coniveau should reduce geometry to residue fields, while transfer might remove finite arithmetic extensions. These routes require careful degree shifts: negative \emph{cohomological} degree here corresponds to positive algebraic $K$-groups. A third route, constructing the Tate $t$-structure and then reading off vanishing, would require precisely the conjectural input we are trying to establish and is not used.

\paragraph{Attempt.}
Write $H^{n,i}(F)$ for rational motivic cohomology of a field, with the comparison to the rational gamma-graded $K$-groups in the question. We work with the standard finite-type smooth schemes of the motivic formalism. For a smooth scheme of finite dimension the coniveau spectral sequence has the form
\begin{equation}\label{eq186:coniveau}
 E_1^{p,q}=\bigoplus_{x\in X^{(p)}}H^{q-p,i-p}(k(x))
 \quad\Longrightarrow\quad H^{p+q,i}(X).
\end{equation}
The shifts arise from codimension-$p$ purity, which subtracts $2p$ from degree and $p$ from weight \eref{186}{1}, \eref{186}{2}. They cannot be replaced by a degree shift of just $p$.

\textbf{Conditional field reduction.} Suppose strong Beilinson--Soul\'e vanishing holds for every field in all positive weights at most $i$. Then $H^{n,i}(X)=0$ for every such smooth $X$ when $n\le0$ and $(n,i)\ne(0,0)$.
For total degree $n=p+q$, the term on the right of \eqref{eq186:coniveau} is
\[
 H^{n-2p,i-p}(k(x)).
\]
If $i-p>0$, it vanishes by the field hypothesis. If $i-p=0$, weight-zero motivic cohomology of a field is zero except in degree zero; the equalities $n-2p=0$, $n\le0$, $p\ge0$ force $n=p=i=0$, the expressly permitted exception. Negative weights vanish in this setting. Thus the entire diagonal of $E_1$ is zero, hence so is its $E_\infty$ diagonal. The coniveau filtration has finitely many steps, so no nonzero infinitely filtered remainder survives. Disconnected finite-type schemes are handled componentwise.

Conversely, an assertion for all smooth schemes over all fields includes $X=\operatorname{Spec}F$. Thus the field problem is not merely a convenient special case: under the usual motivic comparison it contains the obstruction to the geometric statement. If ``smooth scheme'' is interpreted beyond the finite-type convention, passage to that enlarged category needs its own continuity or descent argument; the finite coniveau proof alone does not establish it.

\textbf{A second reduction: algebraic extensions do not kill a rational class.} Let $L/F$ be finite. Restriction and transfer satisfy
\begin{equation}\label{eq186:transfer}
 \operatorname{cor}_{L/F}\operatorname{res}_{L/F}=[L:F]
 \quad\text{on }H^{n,i}(F).
\end{equation}
This is the finite-field-extension degree formula in rational motivic cohomology \eref{186}{1}, Section~4. Therefore restriction is injective. For an algebraic closure $\overline F$, continuity identifies a class vanishing over $\overline F$ as one already vanishing over some finite extension. Equation~\eqref{eq186:transfer} then shows it was zero over $F$. The same reasoning works with a separable closure using finite separable extensions only. Inseparable degrees are invertible with rational coefficients; no characteristic-zero hypothesis on the base is inserted.

Consequently it would suffice to prove the field vanishing over algebraically closed fields of arbitrary characteristic and transcendence degree. This is a reduction, not a claim that algebraically closed fields have trivial higher $K$-theory. In particular, taking an algebraic closure does not remove the transcendental field parameters that the cited paper identifies as difficult.

\textbf{Low-weight consistency check.} Weight zero gives only ordinary degree-zero classes. The standard identification $\Z(1)\simeq\mathbf G_m[-1]$ implies $H^{n,1}(F)=0$ for $n\le0$, and the field reduction proves the same statement for smooth finite-type schemes \eref{186}{2}, Lecture~4. For projective space one can check the indexing independently:
\[
 H^{n,i}(\PP^r_F)=\bigoplus_{j=0}^{r}H^{n-2j,i-j}(F).
\]
No extra $H^{0,0}$ summand occurs at $n\le0$, positive $i$, because it would require $n=2j$, $i=j$ and hence $n>0$. This catches an otherwise tempting sign error in the purity shift.

\paragraph{Stress test.}
Regularity gives vanishing of negative algebraic $K$-groups, but the groups here involve $K_{2i-n}$, whose index is positive for $i>0$, $n\le0$. Adams operations splitting a positive $K$-group into eigenspaces do not prove any particular eigenspace zero. Similarly, transfer proves injectivity after algebraic extension, not vanishing after extension. The coniveau argument requires a convergent finite-dimensional filtration; it does not invoke an unjustified exchange of an infinite intersection with a limit. The argument proves no positive-weight field vanishing beyond the stated known low-weight case. The stronger coefficient comparisons discussed in the literature are not needed and are not inferred from a bare rational calculation.

\paragraph{Outcome.}
\textbf{Outcome: Conditional reduction.}
The geometric target reduces to the field vanishing, and rational restriction permits a further reduction to algebraically closed fields, with arbitrary transcendence degree still present. The index bookkeeping, weight-one check, and transfer argument are supplied explicitly. The unresolved input is the general positive-weight, nonpositive-degree field vanishing; the standard formalism does not manufacture it. These are structural reductions rather than a claim of a new vanishing theorem.

% END RESEARCH ADDITION ATTEMPT-186
\subsection{Sources}
\begin{itemize}\raggedright
\item[\textnormal{[S1]}]\hypertarget{source:186:1}{} F. Deglise, "Generic motives and motivic cohomology of fields", arXiv:2408.06233 [math.AG], v1 dated 12 August 2024.
\url{https://arxiv.org/pdf/2408.06233}.
{\small\textit{Catalogue formulation source.}}
% BEGIN RESEARCH ADDITION REFERENCES-186
\item[\textnormal{[E1]}]\hypertarget{extra:186:1}{} Fr\'ed\'eric D\'eglise, \emph{Generic motives and motivic cohomology of fields}, arXiv:2408.06233v2, Sections 3--4 (generic motives and functorialities), Section 5.2 (Beilinson--Soul\'e and related vanishing questions).
\url{https://arxiv.org/html/2408.06233v2}.
{\small\textit{Research reference; consulted 27 September 2026.}}
\item[\textnormal{[E2]}]\hypertarget{extra:186:2}{} Carlo Mazza, Vladimir Voevodsky and Charles Weibel, \emph{Lecture Notes on Motivic Cohomology}, Clay Mathematics Monographs 2, American Mathematical Society, 2006. Theorem 4.1 and Corollary 4.2 (weights zero and one); Section 14.5 (Gysin and projective-bundle formalism).
\url{https://www.claymath.org/library/monographs/cmim02.pdf}.
{\small\textit{Research reference; consulted 27 September 2026.}}
% END RESEARCH ADDITION REFERENCES-186
\end{itemize}

\end{document}
